% 第 15 章骨架：导引框已定，正文待写。写作计划见 docs/PLAN.md“第 15 章”，写法见 docs/STYLE.md。
\chapter{早期金融革命：荷兰、英国与可信的国债}\label{ch:finrevolution}

\begin{chapterguide}
\guidequestion{荷兰和英国怎样学会以低利率借到大量的钱？}
\begin{guidegoals}
\item 描述阿姆斯特丹银行（1609）的银行货币与荷兰东印度公司（1602）的可转让股份；
\item 用光荣革命后的制度变化（North \& Weingast 1989）解释英国借款成本的下降，并介绍相关争论；
\item 说明英格兰银行（1694）与长期国债的作用；
\item 比较荷兰、英国与同时期西班牙、法国的公共财政。
\end{guidegoals}
\guideprereq{\chref{ch:silver}\enspace\nameref{ch:silver}；\chref{ch:prices}\enspace\nameref{ch:prices}}
\end{chapterguide}

\chapterstub
